% Chapter 30: Jacob Tsimerman — Shimura Varieties and Arithmetic Geometry
\let\LocalMacro\relax

\chapter{Jacob Tsimerman: Shimura Varieties and Arithmetic Geometry}

\section{Prerequisites}
This chapter assumes familiarity with:
\begin{itemize}
\item Algebraic geometry (varieties, abelian varieties)
\item Number theory (Galois representations, modular forms)
\item Model theory (o-minimal structures, Zilber's principle) --- helpful but not essential
\end{itemize}

\section{Overview}
Jacob Tsimerman (born 1988, Russian-born Canadian) was recognized for proving one of the central statements in the theory of Shimura varieties and for fundamental contributions to Hodge theory. His work unites three fields---model theory, Diophantine geometry, and the theory of Shimura varieties---that had previously developed largely independently.

Tsimerman's favorite mathematical metaphor captures the spirit of his work: \emph{if Shimura varieties could vibrate like physical objects, each of their resonant frequencies would correspond to a specific Diophantine equation.} Proving theorems about Shimura varieties thus amounts to understanding the ``spectrum'' of Diophantine equations.

\section{Shimura Varieties}

Shimura varieties are higher-dimensional geometric objects that parametrize abelian varieties with additional structure (such as a polarization, level structure, and endomorphism data). They play a central role in the Langlands program, serving as the geometric stage on which automorphic forms and Galois representations interact.

For example, modular curves (the simplest Shimura varieties) parametrize elliptic curves with level structure. Higher-dimensional Shimura varieties parametrize abelian varieties with more complex structures---they are, in a sense, ``higher-dimensional modular forms.''

\section{The André--Oort Conjecture}

The André--Oort conjecture concerns the structure of \emph{special points} on Shimura varieties---points whose corresponding abelian varieties have unusually large endomorphism algebras (analogous to how the equilateral triangle is ``extra-symmetric'' among triangles).

\begin{conjecture}[André--Oort, 1989]
If $Z$ is a closed subvariety of a Shimura variety $S$, and $Z$ contains a dense set of special points, then $Z$ is itself a special subvariety (i.e., a Shimura subvariety).
\end{conjecture}

The conjecture says that special subvarieties are exactly the ``algebraically defined'' subvarieties of Shimura varieties. It was one of the central open problems in arithmetic geometry for over three decades.

\begin{historicalbox}
The André--Oort conjecture was formulated independently by Yves André (1989) and Frans Oort (1989). It connects the theory of Shimura varieties to transcendence theory and model theory in a fundamental way.
\end{historicalbox}

\subsection{The Proof}

\begin{theorem}[Pila--Shankar--Tsimerman, 2021 \cite{pilatsimerman2021}]
The André--Oort conjecture holds for all Shimura varieties.
\end{theorem}

The proof combined three major ingredients:

\begin{enumerate}
\item \textbf{Pila--Zannier method}: An approach from Diophantine geometry that reduces algebraic statements to counting rational points on definable sets. The key idea is to show that a hypothetical counterexample would contain too many rational points.
\item \textbf{O-minimality}: The theory of o-minimal structures (from model theory) controls the complexity of definable sets, enabling effective counting of rational points via the Pila--Wilkie theorem.
\item \textbf{Zilber's principle}: A model-theoretic axiom asserting that sufficiently saturated geometric structures behave like algebraically closed fields. This allows one to transfer results from the o-minimal setting to the algebraic setting.
\end{enumerate}

\begin{highlightbox}
The proof of the André--Oort conjecture was a tour de force that brought together model theory, Diophantine geometry, and the theory of Shimura varieties---three fields that had previously developed largely independently.
\end{highlightbox}

\section{Hodge Theory and the Griffiths Conjecture}

Tsimerman also proved a major conjecture of Phillip Griffiths concerning the relationship between the \emph{Griffiths group} (a measure of how far algebraic cycles are from being algebraically equivalent to each other) and the \emph{Abel--Jacobi map}.

\begin{theorem}[Bakker--Brunebarbe--Tsimerman, 2023 \cite{bakker2023}]
The Griffiths group of a generic projective hypersurface of sufficiently high degree is torsion-free.
\end{theorem}

This result has direct implications for the Hodge conjecture---one of the seven Millennium Prize Problems worth \$1 million---by constraining the structure of algebraic cycles on algebraic varieties.

\section{Impact}
\begin{itemize}
\item \textbf{Arithmetic geometry}: The André--Oort proof established a new paradigm for studying special subvarieties.
\item \textbf{Model theory}: The use of Zilber's principle and o-minimality demonstrated the power of model-theoretic techniques in algebraic geometry.
\item \textbf{Langlands program}: Tsimerman's work provides tools for studying Galois representations attached to automorphic forms.
\item \textbf{Hodge theory}: The Griffiths result constrains the structure of algebraic cycles, one of the central objects in algebraic geometry.
\end{itemize}

\section{Exercises}

\begin{exercise}[Basic]
What is a Shimura variety? Explain how the modular curve $Y(N)$ (parametrizing elliptic curves with level-$N$ structure) is the simplest example.
\end{exercise}

\begin{exercise}[Intermediate]
Explain the Pila--Zannier method: how does counting rational points on definable sets help prove algebraic statements about special subvarieties?
\end{exercise}

\begin{exercise}[Challenge]
The André--Oort conjecture classifies special subvarieties of Shimura varieties. Explain how this connects to the Langlands program and to the proof of Fermat's Last Theorem.
\end{exercise}

\section{Timeline}

\begin{tabularx}{\linewidth}{lX}
\toprule
\textbf{Year} & \textbf{Event} \\ \midrule
1989 & André and Oort formulate the André--Oort conjecture \\
2021 & Pila, Shankar, and Tsimerman prove the André--Oort conjecture \\
2023 & Bakker, Brunebarbe, and Tsimerman resolve the Griffiths conjecture \\
2026 & Jacob Tsimerman awarded the Fields Medal at ICM Philadelphia \\
\bottomrule
\end{tabularx}

\section{Citations}

\begin{enumerate}
\item Pila, J., Shankar, A. N., \& Tsimerman, J. (2021). ``Ax--Schanuel for Shimura varieties.'' arXiv:2109.08788.
\item Bakker, B., Brunebarbe, Y., \& Tsimerman, J. (2023). ``Ultra-log-rigid spaces and subvarieties of generic projective hypersurfaces.'' Inventiones Mathematicae, 232, 163--228.
\item André, Y. (1989). ``G--functions and geometry.'' Aspects of Mathematics, Vieweg.
\item Oort, F. (1989). ``Special points on Shimura varieties.'' (Unpublished manuscript).
\item Pila, J., \& Wilkie, A. J. (2006). ``The rational points of a definable set.'' Duke Mathematical Journal, 133(3), 591--616.
\end{enumerate}
